\infographicchapter{B-Scan Ultrasonography}{bscan.png}
\label{ch:infographic-bscan}

\section{Overview}

Eye B-scan ultrasonography---short for brightness-scan ultrasound, and sometimes listed on orders as ocular or orbital ultrasonography---is a gentle, painless imaging test that creates two-dimensional cross-sectional pictures of the inside of your eye and the area right behind it. A specialist places a small amount of cool gel on your closed eyelid and rests a smooth handheld wand against it. The wand sends harmless, high-frequency sound waves into the eye and listens to the echoes that bounce back, turning them into a black-and-white picture where brighter spots show stronger echoes.

Because these sound waves travel about 4--5 cm deep, B-scan ultrasound gives doctors a clear view of the back of the eye---including the clear gel inside the eye (the vitreous), the light-sensing layer at the back (the retina), the blood-vessel layer behind it (the choroid), the tough white outer wall (the sclera), and the optic nerve. It is especially helpful when something blocks the doctor's normal view with a light, such as a dense cataract (clouded lens), corneal scarring, or bleeding inside the eye. Doctors commonly use it to check for a detached retina, a shifting vitreous gel, bleeding inside the eye, eye tumors, hidden foreign objects, and problems behind the eyeball, as well as to plan eye surgery.

Its biggest strengths are that it is fast, painless, widely available, radiation-free, and shows live movement inside the eye. Its main limits are that picture quality depends on the skill of the person holding the wand, it shows less microscopic detail than optical coherence tomography (OCT) for the center of the retina, and dense objects can cast shadows on the image. Most importantly, doctors must not press the wand directly on an eye that may have an open wound or rupture, since pressure could push out tissue from inside the eye.

\section{Underlying Physical Principles}
\begin{itemize}
    \item \textbf{Sending and Listening to Sound Waves:} Inside the handheld wand sits a tiny crystal. When a gentle electrical pulse hits this crystal, it vibrates and sends out brief, high-pitched sound waves (8--12 MHz) that are far too high for human ears to hear. When echoes bounce back from inside the eye, the same crystal catches them and turns them back into electrical signals.
    \item \textbf{Why Tissues Create Echoes:} Sound waves bounce back whenever they hit a boundary between two different kinds of tissue---for example, where clear eye gel meets the retina or the outer wall of the eye. Smooth surfaces reflect strong, crisp echoes, while cloudy areas like bleeding or tumors scatter softer echoes in many directions.
    \item \textbf{Balancing Depth and Detail:} Higher-frequency sound waves show finer detail, while lower-frequency waves travel deeper. Using 8--12 MHz strikes the right balance: it reaches 4--5 cm deep to show the back of the eye and the front of the eye socket while still showing fine structures clearly.
    \item \textbf{Measuring Distance by Echo Timing:} The machine calculates how deep each structure sits by timing how long it takes an echo to travel into the eye and bounce back. Because sound travels at a steady speed in eye tissues (about 1,540 meters per second), longer travel times mean deeper structures.
    \item \textbf{Turning Echoes Into a Live Picture:} The scanner boosts faint echoes from deeper tissues so the whole image looks even, then translates echo strength into shades of gray---black for clear fluid with no echoes, and bright white for dense, reflective walls.
\end{itemize}

\section{Imaging Mechanism}
\begin{enumerate}
    \item \textbf{Applying Gel and Placing the Wand:} The examiner puts a drop of cool, sterile water-based gel on your closed eyelid (or uses a small fluid-filled eye cup). The gel removes tiny air pockets so sound waves pass smoothly into the eye.
    \item \textbf{Sweeping the Sound Beam:} A tiny vibrating crystal inside the wand tip sweeps rapid sound pulses in a fan shape across the eye while the examiner gently tilts and turns the wand to view every angle.
    \item \textbf{Catching and Boosting the Echoes:} Between pulses, the wand listens for returning echoes. The scanner amplifies those faint signals, balances the brightness from front to back, and converts them into digital data.
    \item \textbf{Building the Cross-Sectional Image:} Built-in software organizes the echo data into a crisp two-dimensional grayscale image, adjusting contrast so both bright outer walls and faint floating opacities inside the eye can be seen clearly.
    \item \textbf{Live Viewing and Measuring:} The finished picture streams live on the screen at 20--60 frames per second. This lets your doctor watch how structures move when you shift your eyes, fine-tune brightness, and measure any spots or growths on the spot.
\end{enumerate}

\section{Key Technical Parameters}
\noindent
\begin{tabularx}{\textwidth}{@{}l X@{}}
\toprule
\textbf{Parameter} & \textbf{Typical Value / Description} \\
\midrule
Physical Principle & High-frequency sound waves bounce off eye tissues to create a live black-and-white picture \\
Energy / Signal Source & Small handheld eye ultrasound wand \\
Typical Frequency / Energy & 8--12 MHz (commonly 10 MHz) for viewing the back of the eye \\
Spatial Resolution & About 0.1--0.2 mm front-to-back and 0.3--0.5 mm side-to-side (reaches 4--5 cm deep) \\
Typical Scan Duration & 5--15 minutes for one or both eyes \\
Contrast Agent & None (only cool, sterile water-based gel on the closed eyelid) \\
Ionizing Radiation & None \\
\bottomrule
\end{tabularx}

\section{Common Symptoms That Prompt This Scan}
\begin{itemize}
    \item \textbf{Sudden Flashes of Light, New Floaters, or a Dark Curtain Over Vision:} Urgent symptoms of a possible retinal detachment or bleeding inside the eye---especially when dense blood or a cataract blocks the doctor's direct view.
    \item \textbf{Severe Vision Loss With a Cloudy Eye:} Painless or painful loss of sight behind a very dense cataract, scarred cornea, or swollen eyelid where light cameras cannot see to the back of the eye.
    \item \textbf{Eye Injury or Suspected Foreign Object:} Pain, swelling, or blurred vision after blunt trauma or a flying metal or glass splinter hitting the eye.
    \item \textbf{Suspected Eye Mass or Bulging Eye:} Checking a suspected tumor behind the retina or unexplained swelling at the back of the eye socket.
\end{itemize}

\section{Clinical Usages}

\subsection{Primary Usages}
\begin{itemize}
    \item \textbf{Seeing Past Cloudy Areas in the Eye:} Checks the vitreous gel, retina, choroid, and optic nerve when a dense cataract, scarred cornea, or bleeding blocks the doctor's normal view with a light.
    \item \textbf{Checking for Retinal and Vitreous Detachment:} Finds a detached retina (where the light-sensing layer pulls away from the back wall) and tells it apart from a harmless age-related shift in the vitreous gel or a deeper fluid pocket, using live eye movement to help confirm the difference.
    \item \textbf{Evaluating Bleeding and Floaters:} Measures how thick and widespread bleeding inside the eye is, looks for hidden retinal tears behind the blood, and spots debris from severe eye infections or harmless calcium-like floaters.
    \item \textbf{Checking Eye Tumors:} Finds and measures growths inside the eye---such as choroidal melanoma, benign blood-vessel growths, tumors that spread from elsewhere, or childhood retinoblastoma---and checks their inner texture.
    \item \textbf{Checking Eye Injuries and Foreign Objects:} Locates wood, glass, plastic, or metal fragments inside the eye and checks for internal damage (only after making sure the eyeball is not torn open, or using an extra-gentle technique if needed).
\end{itemize}

\subsection{Secondary Usages}
\begin{itemize}
    \item \textbf{Optic Nerve Checks:} Spots buried calcium deposits in the optic nerve head (optic disc drusen), checks nerve cupping, and measures swelling around the optic nerve sheath that can signal high pressure around the brain.
    \item \textbf{Planning Eye Surgery:} Confirms that the retina is safely attached and maps scar tissue pulling on the retina before retinal surgery or dense cataract removal.
    \item \textbf{Tracking Eye Tumors Over Time:} Measures tumor height and width at follow-up visits to see how well radiation plaques, laser treatment, or other therapies are working.
    \item \textbf{Eye Socket and Outer Wall Conditions:} Checks for deep inflammation at the back of the eyeball (posterior scleritis), swollen eye muscles in thyroid eye disease, and lumps in the front part of the eye socket.
\end{itemize}

\section{Advantages}
\begin{itemize}
    \item \textbf{Sees Through Cloudiness and Bleeding:} Clearly shows the back of the eye even when a dense cataract, corneal scar, or internal bleeding blocks a standard eye exam.
    \item \textbf{Shows Live Movement:} Watches how floating membranes move right after you shift your gaze, which helps doctors tell a harmless vitreous shift apart from a serious retinal detachment.
    \item \textbf{No Radiation and Noninvasive:} Uses gentle sound waves with no X-rays and no contrast dye, making it completely safe for children, pregnant patients, and repeat follow-up visits.
    \item \textbf{Fast, Portable, and Affordable:} Small clinic or bedside machines allow quick answers in an eye office or emergency room at a much lower cost than a CT or MRI scan.
    \item \textbf{Helpful Tissue Clues:} Echo brightness and shadowing patterns give doctors valuable clues to tell apart different eye tumors, calcium deposits, and non-metallic foreign objects.
\end{itemize}

\section{Disadvantages}
\begin{itemize}
    \item \textbf{Depends on the Examiner's Skill:} Getting clear, accurate pictures depends heavily on how carefully the examiner angles the wand, adjusts the brightness, and reads the live video.
    \item \textbf{Less Fine Detail Than Light-Based Scans:} Cannot show the microscopic layers of the central retina (the macula) as sharply as optical coherence tomography (OCT).
    \item \textbf{Echo Shadows and False Lines:} Calcium spots, artificial lens implants, or surgical silicone oil inside the eye can cast shadows or bounce echoes back and forth, hiding areas behind them or mimicking disease.
    \item \textbf{Limited View of the Very Front of the Eye:} Standard 8--12 MHz wands are built for the back of the eye and give less detail of the cornea and iris; a special higher-frequency ultrasound (35--50 MHz) is needed for the front of the eye.
    \item \textbf{Unsafe to Press on an Open Eye Wound:} Pressing a wand against a suspected ruptured eyeball risks pushing out internal eye tissue, so the scan must be skipped or done with extreme care and zero pressure.
\end{itemize}

\section{Contraindications}

\subsection{Absolute}
\begin{itemize}
    \item \textbf{Direct Wand Pressure on an Open Eye Injury:} Pressing an ultrasound wand over a suspected or confirmed open-globe rupture or a freshly unsealed eye surgery wound is unsafe because pressure can push out delicate contents from inside the eye.
\end{itemize}

\subsection{Relative}
\begin{itemize}
    \item \textbf{Severe Eyelid Injury or Active Infection:} Serious eyelid cuts, active skin infection around the eye, or a severe corneal ulcer require skipping direct eye contact; if the scan is truly needed, it is done gently over a closed lid using a sterile wand cover and single-use sterile gel.
    \item \textbf{Allergies to Supplies or Numbing Drops:} If you are allergic to latex wand covers, gel preservatives, or numbing eye drops (used only if scanning directly on the eye), the clinic will switch to latex-free and preservative-free alternatives.
\end{itemize}

\section{Patient Safety Guidelines}
\begin{itemize}
    \item \textbf{Protecting Injured or Recently Operated Eyes:} Examiners avoid direct wand pressure over any suspected open eye wound or fresh surgical incision. When a scan is medically essential, they either defer it or float the wand gently over a thick layer of sterile gel on the closed eyelid.
    \item \textbf{Safe Eye-Specific Power Limits:} Clinics use dedicated eye presets that stay strictly within FDA and AIUM safety limits for eye ultrasound (thermal index at or below 1.0, mechanical index at or below 0.23, and average intensity at or below 50 mW/cm$^2$) and keep the scan as brief as possible.
    \item \textbf{Cleanliness and Infection Prevention:} The wand is thoroughly cleaned and disinfected between every patient following AIUM, CDC, and manufacturer rules. Single-use sterile gel packets and sterile wand covers are used whenever scanning near wounds, infections, or healing surgical eyelids.
    \item \textbf{Keeping the Surface of the Eye Safe:} Routine B-scans are done gently over closed eyelids rather than directly on the eyeball. If direct contact is ever needed, numbing drops are used and pressure is kept very light to avoid scratching the cornea.
    \item \textbf{Allergy and Comfort Checks:} Staff ask about latex or gel sensitivities beforehand and avoid pressing firmly on the eyeball, which could briefly raise eye pressure or slow your heart rate through a natural eye-heart reflex.
\end{itemize}

\section{Prerequisite Tests}
\begin{itemize}
    \item \textbf{No Blood Tests or Lab Work Needed:} Because eye B-scan ultrasound is noninvasive, radiation-free, and uses no contrast dye, you do not need blood tests, kidney checks, or a pregnancy test.
    \item \textbf{Standard Eye Exam First:} Before the ultrasound, your eye doctor usually checks your vision, examines the front of your eye with a slit-lamp microscope, and tries to look at the back of your eye to see what is blocking the view and what question needs answering.
    \item \textbf{Ruling Out an Open Eye Wound After Trauma:} After a serious eye injury or suspected metal splinter, doctors carefully inspect the eye and often do a quick, non-touch orbital CT scan first to make sure the eyeball is intact before any ultrasound wand comes near it.
    \item \textbf{Reviewing Past Eye Surgeries:} Telling your doctor about prior cataract surgery, lens implants, or silicone oil from past retinal surgery helps them recognize harmless echo shadows during the scan.
\end{itemize}

\section{Preparing for the Scan}
\begin{itemize}
    \item \textbf{No Fasting or Medicine Changes:} You can eat, drink, and take your regular pills and eye drops as usual unless your eye doctor tells you otherwise.
    \item \textbf{Remove Contact Lenses and Eye Makeup:} Take out soft or hard contact lenses before the exam (bring your case or glasses), and skip heavy mascara, eyeliner, or thick eyelid creams that can interfere with the smooth gel layer.
    \item \textbf{Share Recent Eye Injuries or Allergies:} Let the examiner know right away if you have had recent eye surgery, an eye injury, a corneal ulcer, or allergies to latex, ultrasound gel, or numbing eye drops.
    \item \textbf{Know What to Expect:} You will sit reclined or lie flat with your eyes gently closed while cool gel is placed on your eyelid. The examiner will ask you to look up, down, left, and right behind your closed lids so they can check every corner of the eye.
    \item \textbf{Driving Home:} A standard closed-lid B-scan does not blur your vision or use sedation, so you can drive yourself home afterward---unless your doctor also gives you dilating eye drops during the same visit.
\end{itemize}

\section{During Scan}
\begin{itemize}
    \item \textbf{Resting Position and Gentle Eye Movements:} You sit reclined or lie back comfortably with both eyes softly closed. Keeping your head still on the headrest, you simply look up, down, left, or right behind your closed eyelids when asked.
    \item \textbf{Relaxed Breathing:} Breathe calmly and naturally. Try not to squeeze your eyelids shut or tense your face, as relaxing your eyelids makes the picture clearer and more comfortable.
    \item \textbf{What You Will Feel:} You will feel cool, smooth gel on your closed eyelid and the gentle glide of the small wand (a little gel may trickle down your cheek, which is harmless and easily wiped off afterward). The test is quiet and painless, with only light touch (or brief mild stinging if numbing drops are used for a rare direct-contact scan).
    \item \textbf{Speak Up Anytime:} Tell the examiner right away if you feel sharp eye pain, too much pressure, lightheadedness, or eyelid irritation so they can lighten their touch or pause.
    \item \textbf{How Long It Takes:} Scanning one or both eyes usually takes just 5--15 minutes, with short pauses as the examiner adjusts settings, watches eye movement, and saves measurements.
\end{itemize}

\section{Normal Appearance}
\begin{itemize}
    \item \textbf{Round Eye Shape and Clear Center:} A healthy eyeball appears as a smooth, round white outline filled with a completely black (echo-free) center representing clear vitreous gel. In older adults, a few faint, harmless floating speckles may appear as the gel naturally ages.
    \item \textbf{Smooth Back Wall:} The retina, choroid, and outer white wall (sclera) lie tightly together as one smooth, curved, bright white line at the back of the eye.
    \item \textbf{Optic Nerve and Eye Socket:} The optic nerve shows up as a neat, dark V-shaped shadow leading backward from the eyeball, surrounded by bright, evenly textured fat and darker eye muscles.
    \item \textbf{Natural Lens:} When viewed from the front, the back curve of the eye's natural lens appears as a faint, smooth white arc near the front of the picture.
\end{itemize}

\section{Observing for Abnormal Patterns}
\begin{itemize}
    \item \textbf{Retinal Detachment vs.\ Vitreous Shift:} A detached retina looks like a bright, thick, taut white ribbon anchored at the optic nerve in the back (often forming a V or funnel shape) that moves very little when the eye shifts. In contrast, a common posterior vitreous detachment looks like a thinner, fainter, wavy line that floats freely and sways easily with eye movement.
    \item \textbf{Choroidal Detachment:} Thick, smooth, dome-shaped bulges curving inward from the sides of the eye---without attaching to the optic nerve and barely moving---point to fluid or blood collecting behind the choroid layer.
    \item \textbf{Eye Growths and Calcium Spots:} A choroidal melanoma often appears as a dome- or mushroom-shaped lump with darker, hollow-looking echoes inside. Childhood retinoblastoma and harmless buried optic nerve calcium deposits (drusen) shine intensely bright even at low brightness settings and cast a dark shadow behind them.
    \item \textbf{Deep Wall Inflammation, Bleeding, and Optic Nerve Pressure:} Fluid collecting right behind the eyeball along the optic nerve creates a dark ``T-sign'' that points to posterior scleritis (inflammation of the back wall). Bleeding inside the eye shows up as tiny floating dots or layered clouds that swirl like a shaken snowglobe when you shift your gaze before settling downward. Measuring the dark sheath around the optic nerve about 3 mm behind the eyeball (wider than roughly 5 mm in adults) can also point to elevated pressure around the brain.
    \item \textbf{Harmless Echo Tricks:} Metal splinters, gas bubbles, and artificial lens implants can create bright repeating streaks (like a comet's tail). Silicone oil left from past retinal surgery slows down sound waves and makes the eyeball look artificially stretched out on the screen.
\end{itemize}

\section{Dealing with Uncertainties}
\begin{itemize}
    \item \textbf{Shadows and Echo Artifacts:} Artificial lens implants, gas bubbles, metal fragments, calcium spots, or surgical silicone oil can cast shadows or distort the apparent length of the eye, hiding parts of the back wall.
    \item \textbf{Dense Bleeding vs.\ Retinal Detachment:} Thick bleeding or clotted gel inside the eye can sometimes look very similar to a shallow retinal detachment, especially if it does not move much or sits near the optic nerve.
    \item \textbf{Small Flat Spots:} Telling a harmless flat freckle (choroidal nevus) apart from a very early eye melanoma---or telling a pocket of blood under the retina apart from a small tumor---can be difficult with standard B-scan alone.
    \item \textbf{Fine-Tuning the Exam:} Examiners clear up borderline views by adjusting brightness up and down, watching live eye movement from several angles, and adding a companion A-scan ultrasound that measures exact echo strength inside a spot.
    \item \textbf{Next Steps When Unsure:} If questions remain, doctors correlate the ultrasound with a direct eye exam once bleeding clears, an OCT scan if light can get through, an orbital CT or MRI scan, or a planned repeat B-scan in 1--4 weeks (for bleeding) or 3--6 months (for small spots).
\end{itemize}

\section{Augmenting Methods}
\begin{itemize}
    \item \textbf{Standardized A-Scan Ultrasound:} Measures the exact strength of echoes inside a lump or membrane, helping doctors tell apart choroidal melanoma, benign blood-vessel growths, and other eye tumors much more accurately.
    \item \textbf{Color Doppler Ultrasound:} Checks for active blood flow inside an eye mass---helping tell a living tumor apart from an old pocket of blood---and checks blood flow in the eye's main arteries.
    \item \textbf{Ultrasound Biomicroscopy (UBM):} Uses much higher-frequency sound waves (35--50 MHz) to take razor-sharp pictures of the very front of the eye, including the iris, drainage angle, and structures right behind the iris.
    \item \textbf{OCT and Dye Angiography:} When the front of the eye is clear enough for light to pass, OCT shows microscopic layers of the central retina, while dye-based eye photographs map leaking or abnormal retinal blood vessels.
    \item \textbf{Orbital CT and MRI:} A CT scan quickly pinpoints metal foreign objects and eye-socket bone fractures without touching the eye, while an MRI gives detailed soft-tissue views of the optic nerve and the area behind the eyeball.
\end{itemize}

\section{Comparison with Alternative Modalities}
\begin{itemize}
    \item \textbf{Versus Optical Coherence Tomography (OCT):} OCT shows far finer microscopic detail of the retina, but light cannot get through dense cataracts or thick bleeding. B-scan is the clear choice whenever the inside of the eye is cloudy or when doctors need to see live movement.
    \item \textbf{Versus Ultrasound Biomicroscopy (UBM):} UBM (35--50 MHz) gives much sharper views of the front of the eye (cornea, iris, and drainage angle), while B-scan (8--12 MHz) reaches deeper to examine the back of the eye and the front of the eye socket.
    \item \textbf{Versus A-Scan Ultrasound:} A-scan uses a single straight sound beam for precise one-dimensional measurements---most commonly measuring the front-to-back length of the eye to calculate artificial lens power before cataract surgery---whereas B-scan sweeps across the eye to create a two-dimensional picture of internal structures and live movement.
    \item \textbf{Versus Orbital CT:} A CT scan is preferred first in severe eye trauma to check for a torn eyeball or metal splinter without touching the eye. B-scan is radiation-free, portable, and much better for watching moving membranes and bleeding inside the eye.
    \item \textbf{Versus Orbital MRI:} MRI shows finer soft-tissue detail behind the eyeball and along the optic nerve, while B-scan is much faster, far less costly, and safe even before metal splinters have been ruled out.
\end{itemize}

\section{Follow-up Steps}
\begin{itemize}
    \item \textbf{Right After the Exam:} The examiner gently wipes the cool gel off your closed eyelid, and you can return to your normal routine right away. If numbing drops were used for a rare direct-contact scan, avoid rubbing your eye until normal feeling returns (usually about 15 minutes) so you do not accidentally scratch the cornea.
    \item \textbf{Urgent Care When Needed:} If the scan finds a retinal detachment, retinal tear, severe internal eye infection, or foreign object, you will be referred the same day to a retinal specialist. Suspected eye tumors are referred promptly to an eye-cancer specialist (ocular oncologist).
    \item \textbf{Follow-Up Tests:} Your doctor may pair the B-scan results with a standard dilated eye exam once cloudiness clears, or order an A-scan ultrasound, OCT, orbital CT, or MRI to confirm details.
    \item \textbf{Routine Monitoring:} If you have bleeding inside the eye, doctors often repeat a B-scan every 1--4 weeks to make sure the retina stays safely attached as the blood clears. Small, stable eye spots or treated tumors are typically rechecked every 3--6 months.
\end{itemize}

\section{Key References}
\begin{itemize}
    \item American Academy of Ophthalmology (AAO). \textit{Preferred Practice Pattern: Posterior Vitreous Detachment, Retinal Breaks, and Lattice Degeneration}.
    \item American Institute of Ultrasound in Medicine (AIUM). \textit{AIUM Practice Parameter for the Performance of an Ultrasound Examination of the Eye and Orbit} and \textit{Guidelines for Cleaning and Preparing External- and Internal-Use Ultrasound Transducers}.
    \item U.S. Food and Drug Administration (FDA). \textit{Marketing Clearance of Diagnostic Ultrasound Systems and Transducers: Ophthalmic Acoustic Output Limits}.
    \item Byrne SF, Green RL. \textit{Ultrasound of the Eye and Orbit}, 2nd ed. Mosby; and StatPearls Publishing: \textit{Ocular Ultrasound (B-Scan)}.
    \item U.S. National Library of Medicine (MedlinePlus). \textit{Eye and orbit ultrasound}. National Institutes of Health (NIH). Available at: \url{https://medlineplus.gov/ency/article/003797.htm}.
\end{itemize}

\section{Further Reading}
\begin{itemize}
  \item Wikipedia contributors. \textit{Medical ultrasound}. Wikipedia, The Free Encyclopedia. Available at: \url{https://en.wikipedia.org/wiki/Medical_ultrasound}.
\end{itemize}
